\section{Evaluation metrics}
\label{sec:metrics}

\begin{lede}
What gets measured at the end of a run, how each number is defined, and which
of them may be compared with which.
\end{lede}

\subsection{Percentile}

Standard linear interpolation, on the sorted values, with
$\mathrm{pos} = (n-1) \cdot q/100$ and $\mathrm{lo} = \lfloor \mathrm{pos} \rfloor$:
\begin{equation}
  \mathrm{percentile}(x, q)
    = x_{\mathrm{lo}} (1 - f) + x_{\mathrm{lo}+1} f,
  \qquad f = \mathrm{pos} - \mathrm{lo}.
\end{equation}

\subsection{Jain's fairness index}

\begin{equation}
  \mathcal{J} = \frac{\left(\sum_i T_i\right)^2}{n \sum_i T_i^2},
  \label{eq:jain}
\end{equation}
computed on each robot's completed-task count \citep{jain1984fairness}.
$\mathcal{J} \in (0, 1]$, equal to 1 exactly when every robot completed the
same number of tasks, and tending to $1/n$ as one robot does everything.

\subsection{Aggregate report fields}

$\Theta = K / \max(1, T)$ (\Cref{eq:throughput}); makespan as the maximum
completion time; mean, median, $p_{95}$ and maximum service time over completed
tasks; total and mean travel distance, incremented once per actual move;
direction switches per 1000 steps; the PIBT counters of
\Cref{sec:pibt-recursion}; and the deadlock counters of
\Cref{sec:deadlock-detection}, with mean recovery time as
$\text{recovery time total} / \text{recovered}$.

Four counters exist specifically to test claims that throughput cannot see, and
are what the hypothesis suite of \Cref{sec:results-hypotheses} scores against:

\begin{center}\small
\begin{tabular}{@{}ll@{}}
\toprule
\textbf{Counter} & \textbf{The claim it tests} \\
\midrule
\code{head\_on\_conflicts} & entry admission reduces head-on meetings (H3) \\
\code{counterflow\_moves} & the measured price of ranking direction rather than enforcing it \\
\code{aisle\_throughput\_per\_1000} & narrow aisles flow better under direction control (H1) \\
\code{starvation\_flips} & flips forced by the maximum green rather than by demand \\
\bottomrule
\end{tabular}
\end{center}

\begin{caveat}
\textbf{Service time is censored.} Every service-time field is computed over
\emph{completed} tasks only, and in the saturated regime of
\Cref{sec:regime} that is a minority of the tasks released. A planner that
completes only the easy ones reports the best mean service time in the table
while delivering nothing. Compare service times only between planners of
comparable throughput; \Cref{fig:censoring} shows what happens otherwise.
\end{caveat}
